\documentclass[11pt]{article}
\usepackage{problem_set_style}
\usepackage{kotex}
\usepackage{booktabs}
\usepackage{array}
\begin{document}
\heading{}{샘플}{}{Spring 2027}{프로젝트 1 · 보고서 샘플 (6주차 제출, 축약본)}

\noindent{\small 아래는 \textbf{가상의 팀}이 쓴 샘플입니다. 수치는 형식을 보이기 위한 예시이며 실제 실험 결과가 아닙니다. 형식과 분량의 기준으로만 참고하세요.}

\noindent\textbf{1. 문제 정의.} 고객 한 명을 행으로, 다음 달 연체 여부를 라벨로 한다. 주 지표는 PR-AUC이다(연체 22\%의 불균형 때문). 베이스라인은 항상 ``연체 아님'' (PR-AUC 0.22).

\noindent\textbf{2. 데이터와 전처리.} 30{,}000행을 60/20/20으로 stratified 분할했다. 표준화의 평균과 표준편차는 \textbf{train에서만} 구해 validation과 test에 적용했다. 범주형(학력, 결혼 상태)은 원-핫 인코딩.

\noindent\textbf{3. 모델 선택과 이유.} 로지스틱회귀는 해석이 쉽고, 특성이 많이 연관되어 있어 L2 정규화 효과를 보이기 좋다. GBDT는 표 데이터에서 강한 비교 대상이다. 거리 기반이 아니므로 GBDT에는 표준화를 쓰지 않았다.

\begin{center}\small
\begin{tabular}{@{}lcc@{}}\toprule
모델 (validation) & PR-AUC & 비고 \\\midrule
베이스라인(항상 0) & 0.22 & \\
로지스틱회귀, 정규화 없음 & 0.46 & \\
로지스틱회귀, L2 ($C=0.1$) & 0.48 & validation으로 선택 \\
GBDT (트리 150개) & 0.55 & validation으로 선택 \\\bottomrule
\end{tabular}
\end{center}

\noindent\textbf{4. 결과.} 최종 test는 \textbf{한 번만} 측정했다: GBDT PR-AUC 0.54, 로지스틱회귀(L2) 0.47. 학습 곡선(그림 1)에서 트리 개수가 150을 넘으면 train은 계속 오르지만 validation은 멈춰 150에서 멈췄다.

\noindent\textbf{5. 수식적 정당화.} L2 정규화가 계수의 크기를 줄여 분산을 낮춘 결과 validation PR-AUC가 0.46에서 0.48로 올랐다(편향은 약간 증가). SHAP으로 한 고객의 예측을 분해하면 최근 두 달의 연체 여부가 가장 크게 기여했다.

\noindent\textbf{6. 재현 방법.} 데이터는 UCI 저장소의 원본 CSV, 시드 42, 분할 60/20/20(stratified), L2 세기 $C=0.1$, 트리 150개. 저장소의 \texttt{run.sh} 한 줄로 표의 숫자가 다시 나온다.

\noindent\textbf{7. 한계와 잘 안 된 부분.} 불균형을 줄이려고 연체 쪽 가중치를 키웠더니 PR-AUC가 오히려 0.02 내려갔다. 가중치가 확률을 왜곡해 순위 지표에는 도움이 안 된다고 진단했다.

\end{document}
